\chapter{78}

\section{Bettmeralp/VS (CH), Donnerstag, 6.
September}

Den Nachmittag verbrachte Louis mit den Schafen. Seine Mitarbeit
erweiterte die Männerrunde zum Quartett. Der freundschaftliche Umgang
untereinander erinnerte ihn an sein Ristorante-Team. Bevor er in
Bitterkeit versank, holte ihn David mit einem kleinen Seitenhieb zurück.
Der löste mit seinen flapsigen Sprüchen wiederholt den einen oder andern
Lacher aus. Vorwiegend bei Mazi. Gelegentlich auch bei Joh, der David
jedoch unmissverständlich an der Leine hielt. An einer langen zwar, wie
Louis beobachtete. Es war offensichtlich, wie viel Respekt David Joh
entgegenbrachte.

Neben dem Versorgen der Tiere und Johs kleineren Trainingseinheiten mit
seinem Hirtenhund, blieben Louis Momente für sich alleine. Joh hatte
Louis mit Inbrunst \emph{Flexas} Fähigkeiten vorgeführt. Bis sie sich an
verschiedenen Orten niederliessen.

Joh und David hatten sich zusammengesetzt. Selbst aus der Entfernung
konnte Louis beobachten, dass sie lebhafte Gespräche führten. David
schien im Wechsel wild zu gestikulieren, entspannte sich wieder und
lachte zwischendurch.

Louis holte seine Trinkflasche aus dem Rucksack. Das Bild der beiden
hatte etwas Anziehendes. Einerseits strahlte Joh fast schon
Väterlichkeit aus. Auch Louis gegenüber. Andererseits drang oft etwas
unbedarft Kindliches durch. Joh versprühte diese Leichtigkeit, die
sofort Kreise zog und jeden rundum anzustecken vermochte.

Wie würde er auf Louis’ wahre Geschichte reagieren? Louis versuchte sich
seinen Gesichtsausdruck vorzustellen, wenn er alles wüsste. Wenn Ciro
platzte und plötzlich Louis vor ihm stehen würde. Stellte sich Johs
Worte und die Konsequenzen vor, die er zu erwarten hatte. Er sah Marco
vor sich. Lügen verwandelten den in eine Furie. Louis traute auch Joh
scharfe Reaktionen zu. Die Gedanken waren beängstigend. Es war zu
gewagt, sich ihm zu offenbaren. Er kannte ihn zu wenig gut. Und warum
überhaupt, sollte Joh ihn verstehen und ihm helfen wollen? Gerade er,
den er getäuscht und für seine eigenen Zwecke hintergangen hatte.
Bewusst missbraucht eigentlich. Louis fühlte sich schuldiger denn je.

Der Donnerstag verlief unaufgeregt. Klar geregelt. Ohne grössere
Störungen. Einzig zu Mittag- und Abendessen trafen sie sich zu viert und
beide Male in der oberen, grösseren Hütte. Es ergab sich ganz von
selbst, dass sie Louis die Koch-Regie überliessen.

Nach dem Nachtessen zogen sich Mazi und David in die untere Hütte
zurück.

Joh und Louis trugen das Geschirr hinein.

«Ich mach die Küche. Könntest du den Brunnen in Ordnung bringen? Du
weisst schon, runterbürsten.»

«Ja, natürlich, mach ich.»

Louis trat vor die Hütte. Er zog sich das T-Shirt über den Kopf, öffnete
die Schnürsenkel seiner Bergschuhe und streifte sie ab. Seine Füsse
sahen geradezu erlöst aus.

Von drinnen hörte er Johs Stimme.

«Ich bring dir gleich die Drahtbürste raus.»

Selbst wenn kaum Wasser in den Brunnen rann, jeder Tropfen war
erfrischend. Louis drehte seine Waden unter den schmalen Strahl.

Joh musste ihm zugesehen haben. Er stand neben Louis und schmunzelte.

Die Programmänderung gab Joh nach dem Nachtessen bekannt.

«Wir steigen schon morgen früh zur \emph{Sunnewaid} runter, Louis. Hier
oben läuft’s,» über den Tisch schaute er Mazi und David an, «ihr schafft
es ohne uns. Unten müssen wir den Schafzaun neu aufstellen. Gut, Louis,
wenn wir das morgen zusammen hinbekommen.»

Die Männer nickten.

Mazi und David verabschiedeten sich. Joh hatte Anrufe zu erledigen.

Die Dusche unter dem freiem Himmel war zum abendlichen Highlight
geworden. Das Tankwasser hatte sich über den Tag ordentlich aufgeheizt.
Der Abend schickte Brise um Brise Frischluft in die Dämmerung. Und wie
Louis dastand, inmitten der gigantischen Berge, die als Ahnung in der
Ferne auf den nächsten Morgen warteten, kam er sich gross vor. Alles
unter ihm schien geradezu unwichtig. Vielleicht schaffte Distanz
Freiheit? Er atmete tief ein. Richtete sich ganz gerade auf, reckte
seinen rechten Arm in die Luft und wünschte sich zur Freiheitsstatue.
Sah, wie er die vergoldete Fackel dem Himmel entgegenstreckte und seine
linke Hand das Buch der Unabhängigkeitserklärung hielt. Überlegte, wie
er es für sich umschreiben könnte. Vielleicht zu einem Rezeptbuch. Um
sich nach genauer Anleitung, Schritt für Schritt, von den vergangenen
Widrigkeiten zu befreien. Mit Rezepten kannte er sich aus. Sie waren
nichts anderes als ausgeklügelte Anleitungen zur Problemlösung. Und er
fragte sich, ob es etwas ausgemacht hätte, wenn sie beide, Giorgia und
er, damals die angedachte USA-Reise tatsächlich gemacht hätten?

Er schüttelte das Wasser ab. Und liess sich lufttrocknen.

Um fünf weckte ihn das Handyrattern. Traumfetzen hingen in der Luft. Sie
zogen vorbei. Und Louis konnte nicht fassen, dass die Nacht schon zu
Ende war. Er musste sich zusammenreissen, um in die Gänge zu kommen.

Den Rucksack vor sich hintragend zirkelte er schliesslich die Leiter ins
Erdgeschoss hinunter. Joh fand er zähneknirschend und in Boxershorts vor
dem Herd stehen. Er versuchte, den Kaffeekocher zusammenzuschrauben.

Mit Seitenblick über den Brillenrand und in Richtung Louis murmelte er
ein \emph{Guten Morgen, Ciro}. Löste die beiden Teile voneinander und
begann von neuem.

«Der Gummiring ist kaputt, hoffe, ich bring das Ding wenigstens noch für
einen letzten Kaffee zusammen.»

Joh presste die Lippen aufeinander. Louis schaute ihm zu.

«Bitte notiere: Gummiringe Kaffeekocher.»

«Mach ich.»

Nach einem kleinen Frühstück im Stehen stiegen sie ab. Louis’ Knie
schmerzte bereits nach den ersten Schritten. Schon wieder. Der Alltag
hier oben verlangte eine ausgeprägt gute körperliche Verfassung. War
eigentlich zum Fitnesstraining ohne Erholungszeit und vor allem bei
Abstiegen zur Herausforderung geworden. Er hatte das Knie vorsorglich
eingebunden. Dennoch kam gelegentlich dieser Schmerz, der sich anfühlte,
als piekten ihn scharfe Messerspitzen von innen. Vor dem bevorstehenden
Weg hatte er Respekt. Joh hatte von steilen Stellen gesprochen.

Heute Abend würde Louis Manuela bitten, sich sein Knie anzusehen.

Auf der \emph{Sunnewaid} angekommen, begegneten sie Manuela und Ava
nur kurz.

Joh umarmte Manuela, als hätte er sie Wochen nicht mehr gesehen. Dann
hielt er sie auf Armeslänge von sich an den Schultern, zog sie erneut an
sich und flüsterte ihr etwas ins Ohr. Sie strahlte und küsste ihn. Ava
schnürte ihre Wanderschuhe.

Ein reichhaltiges Frühstück stand auf dem Tisch. Gedeckt für zwei.

«Also dann, guten Appetit euch beiden. Wir gehen dann mal. Bis später.»

Manuela griff nach einer Einkaufstasche. Sie verabschiedeten sich. Joh
packte sie erneut und strich ihr übers Haar. Die beiden Frauen wollten
sich unverzüglich für Einkäufe nach Brig auf den Weg machen. Nach
einigen Metern drehte sich Manuela um. Sie winkte.

«Wir nehmen ein Taxi zurück.»

Joh nickte. Louis’ Blick schweifte zum abgestellten Autowrack, diesem
alten BMW, über dessen Existenz er sich nach wie vor wunderte. Er hatte
oft daran gedacht, nachzufragen, was das Gefährt hier sollte. Und es
immer wieder vergessen.

«Du könntest die beiden auch mit dem BMW abholen,» Louis grinste, «das
wär doch eine tolle Überraschung, nicht?»

Joh blickte amüsiert zu Louis hinüber.

«Haha,..mit Smoking und Zylinder.»

Sie setzten sich gegenüber an den Tisch vor dem Haus. Joh füllte ihre
Tassen mit Kaffee.

«Wie kommt das alte Auto eigentlich hier rauf?»

Joh stellte die Tasse ab und griff sich ein Stück Brot.

«Hm, der BMW war inklusive.»

«Hä?»

«Damals zur \emph{Sunnewaid}. Den hatte mein Onkel Theo in jungen Jahren
erworben. Wie er den heimlich auf die Bettmeralp hochfahren konnte,
verriet er niemandem. Wenn ich ihn danach fragte, grinste er jedes Mal
süffisant. So war er. Aber noch verwunderlicher ist, dass er ihn nicht
gleich wieder wegbringen musste, weg aus diesem autofreien Gebiet. Der
Gemeindepräsident meinte mal, Onkel Theo habe die richtigen Leute hier
oben gekannt. Und ein bisschen ähnlich geht es auch mir. Der Gemeinderat
ist mir gut gesinnt,» er kratzte sich am Kinn, «obwohl, in regelmässigen
Abständen hält mich der eine oder andere an, den Schrotthaufen doch
endlich zu entsorgen.»

Louis schaute Joh ungeduldig an.

«Und?»

«Ich bring’s fast nicht übers Herz. Der Wagen ist ein Stück Theo, ein
Stück \emph{Sunnewaid}-Geschichte. David will ihn mir immer wieder mal
abluchsen. Er stellt sich vor, ihn fahrtüchtig herzurichten. Aber ich
glaube, David hat keine Vorstellung davon, wie desolat sein Zustand
ist.»

Joh machte eine abwehrende Handbewegung.

«Manchmal ist es ganz gut, abzuwarten, bis sich die Dinge selbst
klären.»

Louis nickte lächelnd. Die Strategie hatte was.

Sie begannen gleich nach dem Frühstück. Die Pfosten waren auf der Weide
bereits platziert. Louis legte sich die Drahtrolle über die Schulter.
Auch jetzt fiel ihm wieder ein, dass Joh ihm die Fortsetzung seiner
Liebesgeschichte mit Manuela in Aussicht gestellt hatte. Immer wieder
verkniff er sich, nachzufragen. Vielleicht würde sich nun eine passende
Gelegenheit ergeben.

Sie starteten die Umzäunung in der Nähe des Hauses.
